\documentclass[10pt,a4paper]{article}
%% ЧЕРНОВИК 3.8, 28.09: у 3 кружок вместо креста; 3.7: 22 задачи (+ мехмат 2.3.1 «ряд и стол»), обязательных 12, ⋆ снята с «пар соседей»; 3.6.1: комментарии владельца к 9–21 (8б, жук без кружка, сноска к 15а); 3.5: комментарии владельца к 1–8 (5б снята, 6 на пункты, 3в и 4 переформулированы); 3.4: листок вчетверо снят (НА БУДУЩЕЕ), правки владельца по атаке; 3.3 — мехмат 2.1.1 как ⋆; 3.2 — замечание Д. Макарова — три части; — версия 2.1 (правки Н. П. Стрелковой) + концепция «пружин»; листок 18α «Вероятность».
%% (listok-derevya/listok-16-alpha.tex). Пока .tex написан руками; в чистовике
%% источником правды станет генератор, как у «Деревьев».
\usepackage{iftex}
\ifPDFTeX
  \usepackage[T2A]{fontenc}
  \usepackage[utf8]{inputenc}
  \usepackage[russian]{babel}
\else
\usepackage{fontspec}
\usepackage{polyglossia}
\setmainlanguage{russian}
\setotherlanguage{english}
\setmainfont{SourceSerifPro}[
  Path           = ../listok-derevya/shrifty/,
  Extension      = .ttf,
  UprightFont    = *_400Regular,
  ItalicFont     = *_400Regular_Italic,
  BoldFont       = *_700Bold,
  BoldItalicFont = *_700Bold_Italic]
\fi
\usepackage[left=11mm,right=13mm,top=18mm,bottom=14mm]{geometry}
\usepackage{amsmath,amssymb}
\ifLuaTeX
  \usepackage[protrusion=true,expansion=true,tracking=true,
              stretch=50,shrink=30,final]{microtype}
\else
  \usepackage[protrusion=true,final]{microtype}
\fi
\emergencystretch=1em
\linespread{0.98}
\frenchspacing
\setlength{\parindent}{0pt}
\pagestyle{empty}

\newlength{\shirinaokna}\setlength{\shirinaokna}{12mm}
\newlength{\vysotaokna}\setlength{\vysotaokna}{0.7\baselineskip}
\newcommand{\okno}{\framebox[\shirinaokna]{\rule{0pt}{\vysotaokna}}}
\newlength{\shirinanomera}\setlength{\shirinanomera}{9mm}
\newlength{\zazor}\setlength{\zazor}{0.6em}
\newcommand{\blok}[1]{\par\noindent #1\par}
\newcommand{\nomer}[2]{\makebox[\shirinanomera][l]{%
  \fbox{\rule{0pt}{\vysotaokna}\textbf{#1}$^{#2}$}}}
\newcommand{\term}[1]{\textit{#1}}
\newcommand{\zvezdy}{%
  \par\medskip
  \centerline{\rule{0.3\linewidth}{0.4pt}\quad $*$\ \ $*$\ \ $*$\quad
              \rule{0.3\linewidth}{0.4pt}}%
  \par\medskip}
%% Задача: номер с маркером, окно, текст.
\newcommand{\zad}[3]{\begin{samepage}\blok{\nomer{#1}{#2}\hspace{\zazor}\okno\hspace{\zazor}#3}\end{samepage}\par\vspace{6pt}}
%% Пункт задачи со своим окном и маркером (как 14 в 16α).
\newcommand{\punkt}[3]{\par\nopagebreak\blok{\okno\hspace{\zazor}\textbf{#1)}$^{#2}$\ #3}}
\newcommand{\tekst}[1]{\blok{#1}\par\vspace{6pt}}

\begin{document}
\centerline{\textit{Школа № 179. Ключики 2026/2027 уч.год}}
\par\medskip\centerline{\rule{0.62\linewidth}{0.4pt}}\par\smallskip
\centerline{\large\textbf{18$\alpha$. \textsc{Вероятность}}}\par\smallskip

\tekst{Бросим симметричную монету тысячу раз. Предсказать, как она упадёт в очередной раз, нельзя, но орлов, скорее всего, будет примерно половина. Бросок монеты~— \term{эксперимент}: процедура, которую можно повторять сколько угодно раз, а результат от раза к разу разный. \term{Событие}~— утверждение о результате эксперимента: после каждого опыта можно сказать, произошло оно или нет: «выпал орёл», «на кубике выпало чётное число». Если в $N$ повторениях эксперимента событие произошло $k$ раз, число $k/N$ называют \term{частотой} события.}

\tekst{У симметричного кубика грани ничем не отличаются друг от друга, поэтому в длинной серии бросков каждая грань обычно выпадает примерно в шестой части случаев. Слова «случайный выбор» («тянут наугад», «садится на случайное место») означают то же самое: в длинной серии опытов каждый вариант выбора встречается примерно одинаково часто.}

\zad{1}{\circ}{На гранях симметричной двадцатигранной кости написаны числа от 1 до 20. Какова примерная частота выпадения \textbf{а)} простого числа; \textbf{б)} числа, делящегося на 3 или на 5?}

\zad{2}{}{Симметричный кубик бросают; если выпала шестёрка, бросок не засчитывают и кубик перебрасывают. Какова примерная частота единицы среди засчитанных бросков? А чётного числа?}

\tekst{Если эксперимент повторять много раз, частота события колеблется около одного и того же числа~— где бы и кем бы эксперимент ни проводился. Это число называют \term{вероятностью} события $A$ и обозначают $P(A)$. Это не теорема, а наблюдение о мире; на нём основана вера в то, что теорию вероятностей можно применять к реальной жизни. Слово «вероятность» используют и в других смыслах: например, говорят о вероятности единичного события, которое нельзя повторить много раз,~— что определённая команда выиграет завтрашний матч. В этом листке мы говорим только о вероятности как о частоте.}

\tekst{Событие «произошло хотя бы одно из событий $A$ и $B$» обозначают $A\cup B$, «произошли и $A$, и $B$»~— $A\cap B$, «$A$ не произошло»~— $\overline{A}$; его называют \term{противоположным} к $A$. События $A$ и $B$ \term{несовместны}, если они не могут произойти одновременно.}

\zad{3}{\circ}{\textbf{а)} Объясните, почему $P(A\cup B)=P(A)+P(B)$, если события $A$ и $B$ несовместны. \textbf{б)} Чему равна вероятность события, которое происходит при любом результате эксперимента? \textbf{в)} Объясните, почему $P(A)+P(\overline{A})=1$.}

\zad{4}{\circ}{Верно ли, что $P(A\cup B)=P(A)+P(B)$ для любых событий $A$ и $B$? Если нет, сформулируйте верное равенство.}

\zvezdy

\zad{5}{\circ}{Две симметричные монеты бросают одновременно. Найдите вероятность того, что выпадут два орла; ровно один орёл.}

\begin{samepage}
\blok{\nomer{6}{}\hspace{\zazor}Бросают два симметричных кубика, красный и синий, и складывают выпавшие числа.}
\punkt{а}{\circ}{Какая сумма имеет большую вероятность: 3 или 10? Во сколько раз?}
\punkt{б}{}{Какая сумма имеет наибольшую вероятность?}
\end{samepage}\par\vspace{6pt}

\zad{7}{\circ}{Десять школьников по очереди тянут билеты из десяти; вытянутый билет на стол не возвращают. Вася выучил один билет. Что ему выгоднее: тянуть первым или последним?}

\begin{samepage}
\blok{\nomer{8}{}\hspace{\zazor}Восемь команд разной силы играют турнир по олимпийской системе: три круга, проигравший выбывает, в каждом матче побеждает более сильная команда. Пары первого круга определяет жеребьёвка. С какой вероятностью}
\punkt{а}{}{в финале встретятся две сильнейшие команды;}
\punkt{б}{}{в финал выйдет третья по силе команда?}
\end{samepage}\par\vspace{6pt}

\begin{samepage}
\blok{\nomer{9}{}\hspace{\zazor}За круглый стол в случайном порядке садятся 10 мальчиков и 10 девочек, среди них Маша. С какой вероятностью}
\punkt{а}{\circ}{справа от Маши сидит мальчик;}
\punkt{б}{}{оба её соседа~— мальчики?}
\end{samepage}\par\vspace{6pt}

\begin{samepage}
\blok{\nomer{10}{}\hspace{\zazor}Жук сидит в вершине проволочного каркаса куба. Каждую минуту он переползает по ребру в одну из трёх соседних вершин, выбирая её случайно. С какой вероятностью}
\punkt{а}{}{через 3 минуты он окажется в вершине, противоположной исходной;}
\punkt{б}{}{через 4 минуты он снова окажется в исходной вершине?}
\end{samepage}\par\vspace{6pt}

\begin{samepage}
\blok{\nomer{11}{}\hspace{\zazor}С какой вероятностью ни разу не выпадут два орла подряд, если монету бросают}
\punkt{а}{}{4 раза;}
\punkt{б}{}{10 раз?}
\end{samepage}\par\vspace{6pt}

\begin{samepage}
\blok{\nomer{12}{}\hspace{\zazor}На выборах бюллетени вынимают из урны по одному в случайном порядке и сразу подсчитывают. С какой вероятностью в какой-то момент подсчёта у кандидатов А и Б окажется поровну голосов, если за А подано}
\punkt{а}{}{3 голоса, а за Б~— 2;}
\punkt{б}{\star}{30 голосов, а за Б~— 20?}
\end{samepage}\par\vspace{6pt}

\begin{samepage}
\blok{\nomer{13}{}\hspace{\zazor}Каждое из чисел 1, 2, \ldots, 8 красят в красный или синий цвет, бросая для каждого числа симметричную монету.}
\punkt{а}{\circ}{Что вероятнее: сумма красных чисел больше суммы синих или меньше?}
\punkt{б}{}{Найдите вероятность того, что сумма красных чисел больше.}
\end{samepage}\par\vspace{6pt}

\zad{14}{\circ}{В урне 3 белых и 2 чёрных шара. Их вынимают наугад по одному, пока не появится белый. С какой вероятностью придётся вынуть не меньше трёх шаров?}

\begin{samepage}
\blok{\nomer{15}{}\hspace{\zazor}Игра состоит из партий; в каждой партии каждая из сторон выигрывает с вероятностью $1/2$, ничьих не бывает. Если игру прервали, приз делят в отношении вероятностей победы, которые были у сторон в момент остановки.}
\punkt{а}{\circ}{Двое играли до трёх выигранных партий и прервали игру при счёте 2:1. Как разделить приз?\footnote{Задачу о разделе приза в прерванной игре обсуждали в 1654 году в переписке Блез Паскаль и Пьер Ферма. Эту переписку считают началом теории вероятностей.}}
\punkt{б}{}{Команды Знатоков и Телезрителей играют до шести побед. Сейчас счёт 2:4 в пользу Телезрителей. С какой вероятностью победят Телезрители?}
\end{samepage}\par\vspace{6pt}

\zad{16}{}{Аня бросает $n+1$ симметричную монету, а Боря~— $n$ монет. С какой вероятностью у Ани выпадет больше орлов, чем у Бори?}

\zvezdy

\tekst{\term{Вероятностная модель} эксперимента~— это конечное множество $\Omega$ и числа $p(\omega)\geqslant 0$, по одному для каждого $\omega\in\Omega$, в сумме равные 1. Элементы $\Omega$~— \term{элементарные исходы}: несовместные события, из которых при каждом проведении эксперимента происходит ровно одно; $p(\omega)$~— их вероятности. Событию отвечает подмножество $A\subset\Omega$, и $P(A)$~— сумма $p(\omega)$ по всем $\omega\in A$. Один эксперимент можно описать разными моделями; если обе правильно описывают эксперимент, вероятность события в них одна и та же.}

\zad{17}{\dag}{Докажите, исходя из определения модели: \textbf{а)} $P(\Omega)=1$ и $P(\varnothing)=0$; \textbf{б)} если $A\subset B$, то $P(A)\leqslant P(B)$; \textbf{в)} $P(\overline{A})=1-P(A)$; \textbf{г)} $P(A\cup B)=P(A)+P(B)-P(A\cap B)$.}

\zad{18}{}{\textbf{а)} Опишите вероятностные модели экспериментов из задач 6 и 10: что такое $\Omega$ и чему равны $p(\omega)$? \textbf{б)} Для задачи 9 постройте две модели: в первой элементарный исход~— рассадка всех двадцати человек, во второй~— пара «кто сидит слева от Маши, кто справа». Сколько исходов в каждой модели? Найдите в обеих ответ на вопрос 9б.}

\zad{19}{\circ}{Бросают две одинаковые монеты. Петя рассуждает так: исходов три~— два орла, две решки, орёл и решка,~— значит, монеты упадут по-разному с вероятностью $1/3$. Вася насчитал четыре исхода и получил $1/2$. Кто прав? В чём ошибка другого?}

\zad{20}{}{Десять человек, среди них Петя и Вася, садятся в случайном порядке \textbf{а)} в ряд; \textbf{б)} за круглый стол. С какой вероятностью между Петей и Васей сидят ровно $k$ человек (за столом~— с одной из сторон)?}

\zad{21}{}{За круглый стол в случайном порядке садятся 4 мальчика и 4 девочки. С какой вероятностью их можно разбить на 4 пары соседей так, чтобы в каждой паре были мальчик и девочка?}

\zad{22}{\star}{Каждый из $n$ пассажиров купил билет на $n$-местный самолёт. Первой вошла рассеянная старушка и села на случайное место. Каждый следующий пассажир садится на своё место, если оно свободно, а если занято~— на случайное свободное. С какой вероятностью последний пассажир сядет на своё место?}

\end{document}
